
\section*{Introduction}

The most crucial and often used analog building element is the operational
transconductance amplifier (OTA). An operational amplifier's output voltage is
calculated by multiplying the input voltage by the gain. On the other hand, the
ratio of the output current to the input voltage is what is known as the gain in
OTA. Because the transconductance amplifier transforms an input voltage into the
current at the output terminal, OTA is a voltage-controlled current source. A
wider bandwidth, a more amazing dynamic range, a greater output resistance, a
wider bandwidth, and no additional phase issues are some of the improved
features that OTAs have over an operational amplifier (OP-AMP). Circuits with
capacitive loads can also employ them \cite{ref1,ref2,ref3}. Because of its
versatility, an OTA is useful for a wide range of electrical systems, such as
filters, oscillators, comparators, and analog-to-digital converters.
Furthermore, on occasion, OTA might act as the core amplifier for the
operational amplifier. Scaled Metal Oxide Semiconductor Field Effect Transistors
(MOSFET)-based OTAs have low gain and excessive power consumption. These OTAs
use a directional cascode architecture to maintain linearity and raise output
impedance without compromising bandwidth. The output resistance of a perfect OTA
is infinite. Iout nearly fully passes through external load capacitance (CL)
because of the very high output impedance node (Rout=Rocasn/Rocasp). We have now
added more cascoded devices and positioned them in the upper and lower branches
in order to obtain a bigger benefit \cite{ref3,ref4,ref5,ref6}.
Conventional OTAs were constructed using several technologies. Its higher power
consumption compared to bipolar-based technology, which generates high-speed
OTA, is a significant problem. The performance of a MOSFET is limited by the
short channel effect, which is more noticeable when the channel length is
shortened to the nanoscale. Alternatively, substantial DC improvements are
achieved with minimal power consumption by using complementary metal oxide
semiconductor (CMOS) technology \cite{ref7,ref8,ref9,ref10}.
However, CNT and the carbon nanotube field-effect device (CNTFET) transistor may
eventually be replaced by one new material and a related tool that has Si and
CMOS. Outstanding mechanical, electrical, and thermal properties suggest a
promising material for the future. It can maintain a current density of more
than 109 A/cm2, has exceptional field emission characteristics, is 100 times
more powerful than steel, and has thermal conductivity comparable to that of a
prospective diamond-based device called the CNTFET
\cite{ref11,ref12,ref13,ref14}. as a supplement to the existing silicon
technology and to make Moore's law more applicable. The CNTFET may be able to
effectively utilize the current CMOS design environment and manufacturing
process scheme since its operating system and device structure are similar to
those of CMOS devices. However, CNTFET implementation is challenging. It is
challenging to grow pure single-wall CNTs; two main challenges are chirality
control and CNT placement. To demonstrate the advantages of CNTFETs, several
researchers have implemented a variety of CNT-based digital modules at the
circuit level
\cite{ref14,ref15,ref16,ref17,ref18,ref19,ref20,ref21,ref22,ref23}. But not much
research has been done in the analog domain, particularly in relation to OTAs.
This research demonstrates the advantages of using CNTFETs for OTA design and
simulation. In this study, CNT-based OTAs at the 32 nm technology node were
developed and simulated using HSPICE. Very accurate Stanford Verilog-A models
have been used to analyze the performance of CNTFET transistors. MOSFET at the
technology node of 32 nm The simulation analysis additionally makes use of the
Stanford Verilog-A CNTFET model and the Berkeley Predictive Technology model
(BSIMv4.6.1) \cite{ref18,ref24,ref25,ref26,ref27,ref28,ref29}.
A PCNTFET-NMOS-TC-OTA with P-CNTFET sources and NMOS sinks, an
NCNTFET-PMOS-TC-OTA with N-CNTFET sinks and PMOS sources, and a pure CNT-OTA
with P-CNTFET sources and N-CNTFET sinks were the three CNT-based TC-OTAs that
were designed and assessed. Pure CNT-TC-OTA performs noticeably better than bulk
CMOS, according to simulation data, whereas hybrid TC-OTAs perform in between.
Pitch (S), diameter (DCNT), and CNT count (N) all have a significant impact on
OTA behavior; the best outcomes occur when N = 20, DCNT = 1.5 nm, and S = 20 nm.
Pure CNT-TC-OTA provide good stable operation, according to frequency response
\cite{ref30}.
There are four parts in the paper. The foundation for the planned work is
outlined forth in Section 1, which provides an overview of CNT and CNTFET
technologies. The operational transconductance amplifier (OTA) and its design
principles are covered in Section 2. For all CNT-based OTA architectures,
Section 3 offers comprehensive simulation studies and findings, as well as
comparisons with conventional CMOS-OTA. The main conclusions are outlined in
Section 4, which also emphasizes the benefits of CNT-based solutions.
